% Pursuing Stacks (cotes 134-2 to 134-8): archive cote 134-4, site page 42 (PDF page 43), typescript page 130.
% typescript. Draft compared: ps3.tex:2109-2183.
% AI-transcribed by Claude Opus 5.5 (claude-opus-5-5, adaptive thinking, effort high), Message Batches API, batch ps2 (msgbatch_01NVwr1Yuo8qJf5rykeGBXPK), 2026-11-12; not yet checked by eye.
\scanpage{134-4/42}
\mspage{130}

53.\quad Before pursuing yesterday's reflection about the $\hand{\pi}_o$-functor and
it's relation to homotopy structures on M, \ it seems more convenient to
interpolate \xout{\unsure{a reflection}} some more or less obvious ``functorialities'' on
the homotopy notions just developped. \ They all seem to turn around the
relationship of such notions in M with a more or less arbitrary functor
\[ u : M \hand{\longrightarrow} N , \]
for the time being I am not making any special assumption on \hand{$u$}. \ In terms
of the four ways we got for describing a homotopy structure in M, \ we get
four corresponding natural conditions of ``compatibility'' of \hand{$u$} with a
given homotopy structure h, \ namely :

(i) \ If \ $f \underset{h}{\hand{\sim}} g$ \ in M, \inserted[typed]{i.e. $(f;g)\hand{\in} R_{\hand{h}}$} \ then \ u(f) = u(g) .

(i') \ If $\underline{I}$ =(I,$\hand{\delta}_o$,$\hand{\delta}_1$)$\hand{\in}\hand{\Sigma}_h$ is a weak homotopy interval, \inserted{then} u($\hand{\delta}_o$) = u($\hand{\delta}_1$).

(ii) \ If f$\hand{\in}$ W$_h$ i.e. f is a homotopism, \ u(f) is an isomorphism.

(ii') If \ X$\hand{\in}$ C$_h$ is a contractible object, \ u(X) is a final object (NB
here we assume that u(e$_M$) \unsure{is} a final object in N).

\inserted{(NB.}It is understood implicitly, \ whenever dealing with intervals and with
contractible \xout{\unsure{elemen}} objects, \ that M admits finite products.\hand{)}

The conditions (i) and (ii) are clearly equivalent, \ and equivalent
to the requitement that u factors into
\hand{(1)}
\[ M\hand{\longrightarrow} M_h \hand{\longrightarrow} N , \]
where M$\hand{\rightarrow}$M$_h$ is the canonical functor of M into the corresponding \uline{homotopy}
\uline{-types} category. \ We also have the tautological implications \ (i)$\hand{\Rightarrow}$(i')
and (ii)$\hand{\Rightarrow}$(ii'). \ Moreover we have \inserted{(trivially)}(ii')$\hand{\Rightarrow}$(i') \ whenever the homotopy
structure on M is a contractibility structure, \ and (i\unsure{'})$\hand{\Longrightarrow}$(i) \ whenever
the homotopy structure on M is a weak homotopy interval structure, \ and
\uline{moreover} \ the functor u commutes to \xout{finite}\inserted[typed]{finite} products. \ All these implica-%
tion are summarized in the diagram
\margin{$*$}
\diagram{\begin{tikzcd}[arrows=Rightarrow]
(i) \ar[r,Leftrightarrow] \ar[d] & (ii) \ar[d] \\
(i') \ar[u, bend left, "\hand{$\mathrm{Hom}_3$}" ', "\hand{$*$}"] & (ii') \ar[l, "\hand{$\mathrm{Hom}_4$}" '] 
\end{tikzcd}\ \unsure{$*$}\hand{,}}

where the symbol Hom$_{\hand{3}}$ or Hom$_4$ \ indicates that the implication qualified by
it is valid provided we assume that the homotopy structure \unsure{on|in} M is in
Hom$_{\hand{3}}$ \ (namely is defined in terms of weak homotopy intervals) \ resp. in
Hom$_4$ (namely is a contractibility structure), \ and where (*) denotes the
extra assumption of commutation of u with finite products.

\margin{42}
